\section*{अध्याय \ref{ch:b3:microbiomes} --- माइक्रोबायोम और सहजीविताएँ}

\begin{solution}{exo:b3:microbiomes:1}
\omterm{def:b3:microbiomes:symbiosis}{सहोपकारिता}: दोनों को लाभ --- शिंबी और राइजोबियम, प्रवाल और शैवाल।
\omterm{def:b3:microbiomes:symbiosis}{सहभोजिता}: एक को लाभ, दूसरे पर असर नहीं --- आँत के अधिकांश जीवाणु, जो
उसी पर पलते हैं जिसे परपोषी पचा नहीं सकता। परजीविता: एक को दूसरे की
कीमत पर लाभ --- नर भ्रूण मारता \emph{Wolbachia}। वही जोड़ी खिसक सकती
है: दीवार पार करता कोई आँत-सहभोजी रोगजनक बन जाता है, स्थिरीकरण छोड़
देने वाला कोई राइजोबियम ग्रंथिका का परजीवी, और ताप-दबाव में पड़ा कोई
शैवाल उसी प्रवाल को हानि करता है जिसने उसे बसाया था।
\end{solution}

\begin{solution}{exo:b3:microbiomes:2}
$H = -(0.5\ln 0.5 + 0.3\ln 0.3 + 0.1\ln 0.1 + 2\times 0.05\ln 0.05) =
0.347 + 0.361 + 0.230 + 0.300 = 1.24$; $e^{H} = 3.4$ प्रभावी
जातियाँ। $D = 0.25 + 0.09 + 0.01 + 0.0025 + 0.0025 = 0.355$; $1/D =
2.8$।
\end{solution}

\begin{solution}{exo:b3:microbiomes:3}
16S सर्वेक्षण बताता है कि वहाँ कौन है, लगभग वंश तक, और उनकी सापेक्ष
वाचन-गिनतियाँ; वह उनके ढोए जीनों के बारे में कुछ नहीं कहता, कवक और
\omterm{def:b3:virology:virus}{विषाणु} चूक जाता है (दूसरे प्राइमर), और प्रतिलिपि-संख्या तथा प्राइमरों
से पक्षपाती है। शॉटगन मेटाजीनोमिकी जीन, प्रकार्य और, पर्याप्त गहराई पर,
पूरे जीनोम तथा प्रभेद बताती है; उसकी प्रति-प्रतिदर्श लागत अधिक है, ऊतक
प्रतिदर्शों में वह परपोषी DNA से भर जाती है, और वह फिर भी नहीं बता
सकती कि कौन-से जीन अभिव्यक्त हो रहे हैं या कौन-सी कोशिकाएँ जीवित हैं।
\end{solution}

\begin{solution}{exo:b3:microbiomes:4}
रेशे का किण्वन उन \omterm{def:b3:microbiomes:gut}{लघु-शृंखला वसा अम्लों} में जो बृहदांत्र और परपोषी को
खिलाते हैं; विटामिन~K और B विटामिनों का संश्लेषण; पित्त अम्लों और
औषधों का रूपांतरण; रोगजनकों के विरुद्ध \omterm{def:b3:microbiomes:gut}{उपनिवेशन-प्रतिरोध};
प्रतिरक्षा-तंत्र की शिक्षा। बिगाड़: \omterm{def:b3:bacteriology:antibiotics}{प्रतिजैविकों} के बाद
\emph{Clostridioides difficile} बृहदांत्रशोथ (शोथकारी आंत्र रोग और
मोटापा भी इसमें फँसाए गए हैं)।
\end{solution}

\begin{solution}{exo:b3:microbiomes:5}
$\binom{20}{5} = \num{15504}$। $10$ वाली जाति: $1 - \binom{10}{5}/
\num{15504} = 1 - 252/\num{15504} = 0.984$; $6$ वाली: $1 - 2002/\num{15504}
= 0.871$; $3$ वाली: $1 - 6188/\num{15504} = 0.601$; $1$ वाली: $1 -
\num{11628}/\num{15504} = 0.25$। $E[S_{5}] = 2.7$ जातियाँ। आच्छादन: एक
एकाकी, $1 - 1/20 = 0.95$।
\end{solution}

\begin{solution}{exo:b3:microbiomes:6}
$400\times 10^{11} = 4\times 10^{13}$ जीवाणु, लगभग \qty{40}{g}। जीन:
$1000$ जातियाँ $\times$ \num{4000}, अतिव्यापन घटाकर --- कोई
$3\times 10^{6}$ भिन्न जीनों की कोटि, यानी मानव \num{20000} से सौ गुना
से अधिक।
\end{solution}

\begin{solution}{exo:b3:microbiomes:7}
\omterm{def:b3:microbiomes:nodule}{नाइट्रोजनेज़} ऑक्सीजन से नष्ट हो जाती है, पर बैक्टेरॉइडों को प्रति
N$_{2}$ सोलह ATP बनाने के लिए श्वसन करना ही पड़ता है। ग्रंथिका का
वल्कुट कोई विसरण-अवरोध बनाता है जो ऑक्सीजन का प्रवेश सीमित करता है, और
\omterm{def:b3:microbiomes:nodule}{लेग्हीमोग्लोबिन}, मिलीमोलर सांद्रता पर, भीतर आई ऑक्सीजन बाँध लेता है
ताकि मुक्त सांद्रता \qty{10}{nM} के पास रहे जबकि बैक्टेरॉइडों तक फ्लक्स
ऊँचा रहे। गुलाबी रंग ऑक्सीलेग्हीमोग्लोबिन है; किसी हरी ग्रंथिका ने
अपना \omterm{def:b3:microbiomes:nodule}{लेग्हीमोग्लोबिन} क्षरित कर लिया है --- वह जीर्ण है और अब स्थिर
नहीं कर रही।
\end{solution}

\begin{solution}{exo:b3:microbiomes:8}
दोनों विफल: न ग्रंथिकाएँ, न \omterm{def:b3:microbiomes:mycorrhiza}{अर्बुस्कुली कवकमूल}, क्योंकि कैल्सियम-स्पंद
\omterm{def:b3:microbiomes:mycorrhiza}{साझा सहजीविता-पथ} का साझा संकेत है। इसलिए वह पथ ग्रंथिकन से पहले का है
--- वह चालीस करोड़ वर्ष पहले \omterm{def:b3:microbiomes:symbiosis}{कवक-सहजीविता} के लिए बना था और बाद में
शिंबी पौधों ने उसे \omterm{def:b3:microbiomes:nodule}{राइजोबिया} के लिए भर्ती कर लिया।
\end{solution}

\begin{solution}{exo:b3:microbiomes:9}
कीयर्स ने सोयाबीन की कुछ ग्रंथिकाओं को आर्गन और ऑक्सीजन के किसी
वायुमंडल में रखा, जिसमें \omterm{def:b3:microbiomes:nodule}{राइजोबिया} नाइट्रोजन स्थिर नहीं कर सकते थे,
जबकि उसी पौधे की बाकी ग्रंथिकाएँ सामान्य रूप से स्थिर करती रहीं; पौधे
ने अस्थिरीकारी ग्रंथिकाओं की ऑक्सीजन-पूर्ति काट दी और उनके \omterm{def:b3:microbiomes:nodule}{राइजोबिया} ने
लगभग आधा ही जनन किया। ऐसे भेद के बिना वे \omterm{def:b3:microbiomes:nodule}{राइजोबिया}, जो महँगा
स्थिरीकरण छोड़ देते, स्थिरीकारियों से तेज़ जनन करते, समष्टि भर में फैल
जाते, और पौधे अंततः ऐसे जीवाणु बसाए रहते जो कुछ भी न देते: \omterm{def:b3:microbiomes:symbiosis}{सहोपकारिता}
ढह जाती।
\end{solution}

\begin{solution}{exo:b3:microbiomes:10}
$\binom{N-N_{i}}{n}/\binom{N}{n} = \prod_{j=0}^{n-1}(N - N_{i} - j)/(N -
j)$, जिसका हर गुणक $1$ से नीचे है; $n$ से $n + 1$ पर जाने से एक और
गुणक, $1$ से नीचे, गुणा हो जाता है, इसलिए चूकने की प्रायिकता गिरती है
और उपस्थिति की प्रायिकता $n$ के साथ चढ़ती है; $n = N$ पर अंश $\binom{N -
N_{i}}{N}$ शून्य है और हर जाति उपस्थित, जिससे $S$ मिलता है। $N_{i},
n \ll N$ के लिए हर गुणक लगभग $1 - N_{i}/N$ है और गुणनफल लगभग $(1 -
N_{i}/N)^{n} \approx e^{-nN_{i}/N} \approx (1 - n/N)^{N_{i}}$ --- यही
प्रतिस्थापन-सहित प्रतिदर्शन का सन्निकटन है।
\end{solution}

\begin{solution}{exo:b3:microbiomes:11}
जो सहजीवी केवल अंडों में यात्रा करता हो उसके लिए कोई नर मृत-अंत है,
इसलिए ऐसा कोई भी लक्षण जो नरों को मादा बना दे, उन्हें उनकी बहनों के
पक्ष में मार दे, या असंक्रमित मादाओं को अपंग कर दे, उस सहजीवी का संचरण
बढ़ा देता है। कोशिकाद्रव्यी असंगति: संक्रमित नरों के शुक्राणु ऐसे
रूपांतरित होते हैं कि युग्मनज तब तक मर जाता है जब तक अंडा वही
\emph{Wolbachia} न ढोता हो, जो उसे उबार लेता है। इसलिए असंक्रमित
मादाएँ वह संतति खो देती हैं जो वे संक्रमित नरों से धारण करती हैं, जबकि
संक्रमित मादाएँ कुछ नहीं खोतीं; यह लाभ संक्रमण की बारंबारता के साथ
बढ़ता है, इसलिए किसी देहली के ऊपर संक्रमण स्थिरण तक फैल जाता है, भले ही
संक्रमित मादाएँ थोड़े कम अंडे दें --- वह लागत असंक्रमित मादाओं की नष्ट
हुई संतति के अंश से ढँक जाती है।
\end{solution}

\begin{solution}{exo:b3:microbiomes:12}
अंतःसहजीवी हर परपोषी-पीढ़ी पर कुछ ही कोशिकाओं की किसी अड़चन से गुज़रता
है और बाहरवालों के साथ कभी पुनर्संयोजन नहीं करता: अपवाह थोड़े हानिकारक
उत्परिवर्तन स्थिर कर देता है, विलोपन समेत, और बदले जीनों का कोई स्रोत
नहीं (मुलर का दंतचक्र)। जिन जीनों के उत्पाद परपोषी जुटाता है, या जो
स्वतंत्र जीवन के काम आते हैं --- आवरण, गतिशीलता, नियमन, अधिकांश मरम्मत
--- वे खटकते नहीं और पहले जाते हैं; अनुवादन-तंत्र और वे जीन जो परपोषी
की ज़रूरत की चीज़ बनाते हैं (\emph{Buchnera} में अनिवार्य ऐमीनो अम्ल)
सबसे देर तक रखे जाते हैं। स्वतंत्र-जीवी नातेदारों की समष्टियाँ विशाल
हैं, उनमें पुनर्संयोजन और जीन-प्रवाह है, और वरण उनके जीनोम बनाए रखता
है। उलटाव के लिए बाहर से नए जीन चाहिए --- अलगाव में असंभव --- इसलिए
परपोषी इसके बजाय किसी क्षरित सहजीवी को किसी ताज़े से बदल देते हैं, जैसा
कीटों की कई वंश-परंपराओं ने किया है।
\end{solution}

\begin{solution}{pb:b3:microbiomes:1}
\textbf{1.} $400\times 10^{11} = 4\times 10^{13}$ जीवाणु, \qty{40}{g}।
\textbf{2.} \num{4000}-जीन का कोई जीनोम लगभग \qty{4}{Mb} है, \qty{4.4}{fg}:
$4\times 10^{13}\times 4.4\times 10^{-15} = \qty{0.18}{g}$ जीवाणु
DNA, जबकि $3\times 10^{13}\times 6.4\times 10^{-12} = \qty{0.19}{g}$
मानव DNA --- लगभग बराबर।
\textbf{3.} $1000\times 0.7\times 4000 + 0.3\times 4000 \approx 2.8\times
10^{6}$ भिन्न जीन, यानी मानव जीनोम से कोई $140$ गुना।
\textbf{4.} $30\times 10 = \qty{300}{mmol}$; $\times 0.9 = \qty{270}{kJ}$;
आहार का \qty{3}{\%}।
\textbf{5.} नहीं --- \qty{3}{\%} से \qty{30}{\%} की व्याख्या नहीं
होती। रोगाणु-मुक्त चूहे कम कुशलता से अवशोषित भी करते हैं, उनके
आँत-हॉर्मोन, वसा-संचय और ऊष्माजनन बदले हुए हैं, और भूख अलग:
\omterm{def:b3:microbiomes:symbiosis}{माइक्रोबायोम} परपोषी की शरीरक्रिया पर काम करता है, केवल
ईंधन-आपूर्तिकर्ता के रूप में नहीं।
\textbf{6.} विभाजन निकासी से संतुलित: कोई स्थायी समष्टि, जिसकी हर
कोशिका रोज़ बदल जाती है। \qty{99.9}{\%} की सफ़ाई के बाद
$4\times 10^{10}$ बचते हैं और दस द्विगुणन लगभग दस दिनों में संख्याएँ
लौटा देते हैं --- पर बचे हुए और सबसे तेज़ बढ़ने वाले हावी हो जाते हैं,
संघटन बदल जाता है (कोई \omterm{def:b3:microbiomes:gut}{डिस्बायोसिस}), और उस अंतराल में \emph{C.
difficile} जैसा कोई हमलावर जम सकता है।
\textbf{7.} गौण जातियाँ $0.4$ को $177$ में बाँटती हैं: हर एक $p = 0.00226$। $H
= -(0.3\ln 0.3 + 0.2\ln 0.2 + 0.1\ln 0.1 + 0.4\ln 0.00226) = 0.361 +
0.322 + 0.230 + 2.437 = 3.35$; $e^{H} \approx 28$।
\textbf{8.} $D = 0.09 + 0.04 + 0.01 + 177\times (0.00226)^{2} = 0.141$;
$1/D = 7.1$। \omterm{def:b3:microbiomes:diversity}{सिम्पसन सूचकांक} प्रचुरताओं का वर्ग करता है, इसलिए वही तीन
हावी जातियाँ उसे तय करती हैं और दुर्लभ शायद ही गिनती में आती हैं;
शैनन दुर्लभ जातियों को अधिक भार देता है।
\textbf{9.} $1 - 60/2000 = 0.97$: प्रति हज़ार कोई $30$ \omterm{def:b3:genomics:ngs}{वाचन} उन जातियों
के हैं जो अब तक देखी नहीं गईं।
\textbf{10.} प्रतिदर्श का आकार बढ़ने पर समृद्धि चढ़ती है क्योंकि दुर्लभ
जातियाँ प्रकट होती रहती हैं; \num{10000} \omterm{def:b3:genomics:ngs}{वाचन} उन जातियों तक पहुँचते हैं
जो \num{2000} चूक जाते हैं। बड़े प्रतिदर्श को प्रमेय से (या
उप-प्रतिदर्शन से) \num{2000} \omterm{def:b3:genomics:ngs}{वाचनों} तक विरल कीजिए --- यदि वही समुदाय
हो तो उसे लगभग $180$ देना चाहिए --- और बराबर गहराई पर तुलना कीजिए।
\textbf{11.} $N_{i} = 1$: $1 - \binom{1999}{1000}/\binom{2000}{1000} =
1 - (2000 - 1000)/2000 = 0.5$। $N_{i} = 5$: लगभग $1 - 0.5^{5} = 0.97$।
\textbf{12.} $e^{2.0} = 7.4$ प्रभावी जातियाँ; $D \ge 0.5^{2} = 0.25$,
लगभग $0.27$: कोई ऐसा समुदाय जिसमें एक हावी जाति और कुछ दर्जन गौण
जातियाँ हैं --- वही ढही, कम-विविधता वाली आँत जिसमें
\emph{C.~difficile} पैर जमाता है।
\textbf{13.} $150\times 8 = \qty{1200}{kg}$ शर्करा प्रति हेक्टेयर, यानी
प्रकाश-संश्लेषित \qty{12}{t} का \qty{10}{\%}।
\textbf{14.} $150\,000/28 = 5360$ मोल N$_{2}$; $16\times 5360 =
\num{85700}$ मोल ATP।
\textbf{15.} $\num{85700}/30 = 2860$ मोल ग्लूकोस, \qty{514}{kg} ---
चुकाई गई शर्करा का लगभग \qty{43}{\%}। बाकी ग्रंथिकाएँ और बैक्टेरॉइड
बनाता तथा बनाए रखता है, प्रति N$_{2}$ आठ इलेक्ट्रॉन अपचायक के रूप में
जुटाता है, और अमोनिया का स्वांगीकरण करता है।
\textbf{16.} $1200\times 16 = \qty{19200}{MJ}$, \qty{150}{kg} N के लिए:
\qty{128}{MJ/kg}, यानी कारख़ाने के \qty{35}{MJ/kg} से लगभग चार गुना ---
पर धूप में, खेत में, बिना किसी ईंधन के चुकाया गया।
\textbf{17.} स्थिरीकारी $0.9\times 1 = 0.9$; छली $0.1\times 0.5 =
0.05$; छली $0.05/0.95 = 5.3\,\%$ हैं --- एक ही मौसम में आधे।
\textbf{18.} छली $0.1\times 1.2^{10} = 0.62$ बनाम स्थिरीकारी $0.9$:
दस मौसमों के बाद \qty{41}{\%}, और आधिपत्य की ओर चढ़ते हुए। प्रतिबंध वरण
को छली के विरुद्ध मोड़ देते हैं; उनके बिना \omterm{def:b3:microbiomes:symbiosis}{सहोपकारिता} क्षरित होती है।
\textbf{19.} $10^{6}\times \qty{20}{pg} = \qty{20}{\micro g}$ कार्बन
प्रति cm$^{2}$ प्रति दिन; \qty{18}{\micro g} प्रवाल तक।
\textbf{20.} श्वसन \qty{15}{\micro g}: अधिशेष \qty{3}{\micro g}
प्रति cm$^{2}$ प्रति दिन, वृद्धि, कैल्सीभवन की कार्बनिक आधात्री,
श्लेष्मा और युग्मकों के लिए।
\textbf{21.} \qty{5}{\%} शैवाल के साथ, दिन में \qty{0.9}{\micro g}
बनाम श्वसित \qty{15}{\micro g}: दिन में \qty{14}{\micro g} की कमी;
\qty{300}{\micro g} का भंडार लगभग तीन सप्ताह चलता है।
\textbf{22.} $+2$ पर चार सप्ताह अधिक बुरे हैं: शैवालों के
प्रकाश-तंत्रों की क्षति तापमान के साथ रैखिक नहीं, तेज़ी से चढ़ती है।
फिर भी वह माप चेतावनी की तरह काम करता है, क्योंकि विवर्णन संचयी दबाव
का समाकलन है और दोनों रूपरेखाओं का अंतर देहली की अनिश्चितता से कम है;
वह कोई आनुभविक सूचकांक है, जो देखे गए विवर्णन पर अंशांकित है।
\textbf{23.} दबाव में पड़ा कोई प्रवाल पहले से कम प्रचुरता में मौजूद
अधिक ताप-सहिष्णु शैवाल-जातियों की ओर खिसक सकता है, या विवर्णन के बाद
उन्हें जल से पा सकता है; सहिष्णु शैवाल कम कार्बन देते हैं, इसलिए प्रवाल
धीमे बढ़ता है, और जो सहिष्णुता मिलती है वह एक-दो अंश की है ---
प्रक्षेपित तापन से कम --- जबकि तापन की चाल प्रवालों के अपने अनुकूलन से
आगे निकल जाती है।
\textbf{24.} $10^{10}$ cm$^{2}$ $\times$ \qty{20}{\micro g} $=
\qty{200}{kg}$ कार्बन प्रति दिन, यानी वर्ष में लगभग \qty{73}{t}।
\textbf{25.} \omterm{def:b3:microbiomes:symbiosis}{माइक्रोबायोम} के किण्वन से आहार का लगभग \qty{3}{\%};
अनुक्रमण-प्रतिदर्श का \qty{97}{\%} आच्छादन; सेम की फ़सल की नाइट्रोजन के
लिए चुकाया प्रकाश-संश्लेषण का \qty{10}{\%}।
\end{solution}
